Congressional apportionment --- the allocation of 435 House seats among the
50 states based on census population counts --- is the foundational step
of every redistricting cycle.
Article~I, \S2 of the U.S.\ Constitution mandates that representatives
``shall be apportioned among the several States according to their
respective Numbers,'' but it does not specify which mathematical method
Congress must use.
Five methods have been proposed, debated, and in some cases enacted:
Huntington-Hill (equal proportions, current law since 1941),
Webster (major fractions), Adams (smallest divisors),
Jefferson (greatest divisors, used 1790--1840), and Hamilton
(largest remainders).
These methods divide into two families: \emph{divisor methods}
(Huntington-Hill, Webster, Adams, Jefferson), which allocate seats by
rounding a population-to-divisor ratio, and \emph{quota methods}
(Hamilton), which allocate seats by rounding exact proportional quotas.

We develop a unified mathematical framework characterising all five methods
along three dimensions: the rounding rule that determines each state's
seat count, the fairness criterion each method optimises, and susceptibility
to the classical apportionment paradoxes.
The central taxonomic result is that divisor methods are provably
immune to the Alabama paradox and related anomalies while quota methods
are susceptible --- a fact that drove Congress's 1941 switch from Hamilton
to Huntington-Hill.
Using 2020 Census data (apportionment population 331{,}108{,}434;
435 seats), we compute all five methods' seat allocations for all 50 states
and identify the three states where methods diverge.
The \textsc{bisect-apportion} Rust crate currently implements the four divisor
methods with floating-point priority comparisons.  Hamilton remains a
mathematical reference point for paradox analysis, and Census-reference
verification fixtures remain future hardening work.
This paper is the foundational overview for J.1--J.6, which study each
method and the paradox theory in depth.
